% Standalone problem document, generated from the adjacent README.md.
% Compile with XeLaTeX. Mathematical content is maintained in Markdown.
\documentclass[11pt,a4paper]{article}
\usepackage[fontsize=10.5pt]{scrextend}
\usepackage[margin=22mm,headheight=15pt,headsep=9mm,footskip=11mm]{geometry}
\usepackage{fontspec}
\setmainfont{texgyrepagella}[Extension=.otf,UprightFont=*-regular,BoldFont=*-bold,ItalicFont=*-italic,BoldItalicFont=*-bolditalic]
\setsansfont{texgyreheros}[Extension=.otf,UprightFont=*-regular,BoldFont=*-bold,ItalicFont=*-italic,BoldItalicFont=*-bolditalic]
\setmonofont{texgyrecursor}[Extension=.otf,UprightFont=*-regular,BoldFont=*-bold,ItalicFont=*-italic,BoldItalicFont=*-bolditalic,Scale=MatchLowercase]
\usepackage{amsmath,amssymb}
\usepackage{unicode-math}
\setmathfont{texgyrepagella-math.otf}
\usepackage{xcolor,microtype,enumitem,needspace,fancyhdr}
\usepackage{bookmark}
\definecolor{ink}{HTML}{17344B}
\definecolor{accent}{HTML}{197D80}
\definecolor{muted}{HTML}{596773}
\hypersetup{colorlinks=true,linkcolor=accent,urlcolor=accent,pdftitle={IV-03 - Polynomial-size
vertex test for inverse M-matrix
intervals},pdfauthor={Open Problems in Numerical Linear Algebra}}
\urlstyle{same}
\setlength{\parindent}{0pt}
\setlength{\parskip}{5pt plus 1pt minus 1pt}
\linespread{1.04}
\setlength{\emergencystretch}{3em}
\widowpenalty=10000
\clubpenalty=10000
\displaywidowpenalty=10000
\allowdisplaybreaks[1]
\setlist{nosep,leftmargin=1.5em}
\providecommand{\tightlist}{\setlength{\itemsep}{0pt}\setlength{\parskip}{0pt}}
\providecommand{\pandocbounded}[1]{#1}
\pagestyle{fancy}
\fancyhf{}
\fancyhead[L]{\sffamily\footnotesize\color{muted}OPEN PROBLEMS IN NUMERICAL LINEAR ALGEBRA}
\fancyhead[R]{\sffamily\bfseries\footnotesize\color{accent}IV-03}
\fancyfoot[L]{\parbox[b]{0.9\textwidth}{\sffamily\fontsize{8}{10}\selectfont\color{muted}Editorial ratings; a bounded search cannot certify that a problem remains unsolved.\\Literature check: 2026-09-15}}
\fancyfoot[R]{\sffamily\footnotesize\color{muted}\thepage}
\renewcommand{\headrulewidth}{0.3pt}
\renewcommand{\headrule}{\hbox to\headwidth{\color{accent}\leaders\hrule height \headrulewidth\hfill}}
\makeatletter
\renewcommand{\section}{\Needspace{5\baselineskip}\@startsection{section}{1}{0pt}{12pt plus 2pt minus 2pt}{4pt}{\sffamily\bfseries\large\color{ink}}}
\renewcommand{\subsection}{\Needspace{4\baselineskip}\@startsection{subsection}{2}{0pt}{9pt plus 1pt minus 1pt}{3pt}{\sffamily\bfseries\normalsize\color{ink}}}
\makeatother
\setcounter{secnumdepth}{0}
\begin{document}
{\sffamily\small\color{muted}Intervals and absolute value equations\par}
\vspace{3pt}
{\raggedright\hyphenpenalty=10000\exhyphenpenalty=10000\sffamily\bfseries\fontsize{21}{24}\selectfont\color{ink}Polynomial-size
vertex test for inverse M-matrix intervals\par}
\vspace{7pt}
{\sffamily\small\textbf{Difficulty:} challenging\quad\textbf{Importance:} interesting
to specialist\par}
{\sffamily\small\bfseries\color{accent}Status: Lean verified\par}
\vspace{4pt}
{\color{accent}\hrule height 0.8pt}
\vspace{6pt}
\textbf{Rating rationale:} Replacing an exponential interval test by a
quadratic vertex family is challenging; verified inverse-positivity has
specialist importance in interval matrix analysis.

\textbf{Area:} interval linear algebra; structured matrices

\subsection{Independently reviewed resolution -
2026-09-11}\label{independently-reviewed-resolution---2026-09-11}

\textbf{Affirmative resolution.} Theorem 1 proves that every interval
member is inverse-M if and only if the \(n^2\) vertices \(C-D_iRD_j\)
are inverse-M. These are contained in the displayed two-sign family, so
the original \(2n^2\) equivalence follows. The proof covers all real
endpoints, every dimension, zero widths, zero entries and reducible
matrices without assuming regularity.

\textbf{Author:} Matthew J. Colbrook, Department of Applied Mathematics
and Theoretical Physics, University of Cambridge. See the
\href{https://github.com/ajt60gaibb/OpenProblemsInNLA/blob/main/references/colbrook-intervals-2026-09-11/manuscripts/IV-03.pdf}{complete
manuscript},
\href{https://github.com/ajt60gaibb/OpenProblemsInNLA/blob/main/references/colbrook-intervals-2026-09-11/verification/reviews/IV-03-review.md}{independent
agent review} and
\href{https://github.com/ajt60gaibb/OpenProblemsInNLA/blob/main/references/colbrook-intervals-2026-09-11/README.md}{submission
record}. The source archive identifies the drafts as AI-generated;
authorship is recorded at the submitter's request. This is agent
verification, not external human peer review or formal proof-assistant
certification.

The difficulty, importance and rating rationale below are historical
assessments of the original open target. The original statement and
dated audits are preserved.

\subsection{Lean verification --- 15 September
2026}\label{lean-verification-15-september-2026}

The complete original question has a formally verified affirmative
answer: for every positive dimension and every ordered pair of real
entrywise interval endpoints, the \(n^2\) negative-sign vertices
characterize the entire interval consisting of inverse M-matrices. This
proves the full original two-sign criterion. Zero widths, zero entries
and reducible matrices are included; no interval regularity or
nonsingularity assumption is added. The manuscript's arithmetic and
bit-complexity claims are outside this formalization.

All four reviewed statements passed LeanCert kernel trust audits,
standard-axiom checks, and
\href{https://github.com/sidneyholden1/OpenProblemsInNLA/actions/runs/34926260380}{fresh
isolated Linux Comparator/default-kernel verification} at the
\href{https://github.com/sidneyholden1/OpenProblemsInNLA/tree/516ad4a0e85c21c7ef34507db9ab3b68b393bb90/intervals-and-absolute-value-equations/IV-03/lean}{immutable
proof revision}. Two independent statement reviews preceded
implementation and two independent nonauthor final code reviews passed.
Actual rejection and sandbox controls passed.

Formalization: Sidney Holden, with OpenAI Codex assistance. Matthew J.
Colbrook retains mathematical authorship as recorded in the source.
Reviews are AI-agent reviews, not external human peer review or official
Tau Ceti endorsement; no source-author endorsement of the formalization
is asserted.

\href{https://github.com/sidneyholden1/OpenProblemsInNLA/tree/516ad4a0e85c21c7ef34507db9ab3b68b393bb90/intervals-and-absolute-value-equations/IV-03/lean}{Proof
and reviews} ·
\href{https://github.com/sidneyholden1/OpenProblemsInNLA/tree/codex/lean-iv03/docs/lean/verification-2026-09-15/IV-03}{Retained
Linux evidence}.

\newpage

\subsection{Problem statement}\label{problem-statement}

A real square matrix \(M\) is a nonsingular M-matrix if its off-diagonal
entries are nonpositive, \(M\) is invertible, and \(M^{-1}\) is
entrywise nonnegative. An inverse M-matrix is the inverse of such a
matrix.

For every \(n\ge1\), let \(L,U\in\mathbb R^{n\times n}\) satisfy
\(L\le U\) entrywise, and set

\[\mathcal A=\{A:L\le A\le U\},\qquad C=(L+U)/2,\qquad R=(U-L)/2.\]

For \(i=1,\ldots,n\), let \(z^{(i)}\) have entry \(-1\) in position
\(i\) and entry \(1\) elsewhere, and write
\(D_i=\mathop{\mathrm{diag}}\nolimits(z^{(i)})\).

Is the following equivalence true?

\[\begin{gathered}
\bigl[\text{every }A\in\mathcal A\text{ is an inverse M-matrix}\bigr]
\\\Longleftrightarrow\\
\left[\begin{gathered}
C+sD_iRD_j\text{ is an inverse M-matrix}\\
\text{for every }i,j\in\{1,\ldots,n\}\text{ and }s\in\{-1,1\}
\end{gathered}\right].
\end{gathered}\]

All entries in the interval vary independently, and degenerate intervals
are allowed. The proposed test uses at most \(2n^2\) matrices; repeated
test matrices are harmless. This would give a small exact certificate
for a structured property needed in interval inversion and verified
linear-system computation.

\subsection{References}\label{references}

Milan Hladík,
\href{https://doi.org/10.1007/978-3-030-31041-7_16}{\emph{An overview of
polynomially computable characteristics of special interval matrices}},
in \emph{Beyond Traditional Probabilistic Data Processing Techniques},
Springer (2020), pp.~295--310;
\href{https://arxiv.org/pdf/1711.08732}{author preprint,
arXiv:1711.08732v1}, §9, Conjecture 1, p.~11, with the vectors defined
in Theorem 24.

Jürgen Garloff, Doaa Al-Saafin, and Mohammad Adm,
\href{https://reliable-computing.org/reliable-computing-28-pp-056-070.pdf}{\emph{Further
Matrix Classes Possessing the Interval Property}}, Reliable Computing
\textbf{28} (2021), 56--70, p.~64, paragraph immediately before IP 4.4.

\subsection{Status check}\label{status-check}

The 2021 paper explicitly retains this conjecture while proving a
different sufficient vertex family of size \(2^{n-1}\). On 2026-09-10,
searches combining the original title, inverse M-matrices, interval
property, conjecture, and 2025/2026 found no resolution. Hladík's
current publication list and the later interval-property paper were
checked. This is a selected-source status check, not an exhaustive
citation audit.

\subsection{Independent audit ---
2026-09-10}\label{independent-audit-2026-09-10}

The full Hladík preprint, §9, Conjecture 1, agrees with the displayed
two-sign vertex family, including the definition of each \(z^{(i)}\).
Garloff--Al-Saafin--Adm, p.~64, expressly leaves that conjecture
unresolved; its exponential family is a different test, not a resolution
for a subclass here. Independent later searches for the exact conjecture
and matrix class found no resolution.
\end{document}
